% ============================================================================

% Paper: Frequent Losers as Cover Bidders in Public Procurement
% Target: RAND Journal of Economics
% ============================================================================

\section{Empirical Strategy}
\label{sec:empirical}

Our empirical strategy has three tiers, ordered by the strength of causal
claims. \emph{First}, we estimate the structural model of cover-bidder
deployment from Section~\ref{sec:structural}, which identifies the
cover-bid distribution and implied markup directly from the bid-level
data (Section~\ref{sec:structural_estimation}). \emph{Second}, we
provide reduced-form evidence using high-dimensional OLS with item, year,
and purchasing-unit fixed effects, supplemented by a leave-one-out
instrumental variable (Section~\ref{sec:reduced_form}). \emph{Third}, we
validate the FL screen's detection performance against CADE cartel
convictions using ROC analysis (Section~\ref{sec:detection_validation}).
Non-structural corroboration is provided by corrected Bajari--Ye
bid-coordination tests (Section~\ref{sec:bajari_ye}).

% ── 5.1 Structural Estimation ────────────────────────────────────

\subsection{Structural Likelihood Estimation}
\label{sec:structural_estimation}

We estimate the mixture model from Section~\ref{sec:likelihood} in two stages using bid-level data. After restricting to losing bids with valid prices and item-year fixed-effect coverage, the estimation sample comprises 23.2 million genuine (non-FL) losing bids and 189,381 FL losing bids (Table~\ref{tab:structural_params}).

\paragraph{Stage 1: Genuine bid distribution.}
We estimate $(\boldsymbol{\beta}, \sigma_g^2)$ from the 23.2 million non-FL losing bids by maximum likelihood under the log-normal specification~\eqref{eq:genuine_dist}:
\begin{equation}
  (\hat{\boldsymbol{\beta}}, \hat{\sigma}_g^2) =
  \operatorname*{arg\,max}_{\boldsymbol{\beta}, \sigma_g^2}
  \sum_{t=1}^{T} \sum_{j \in \mathcal{G}_t}
  \log \phi\!\left(
    \frac{\log b_{jt} - \mathbf{x}_j'\boldsymbol{\beta}
    - \hat{\alpha}_i - \hat{\lambda}_t}{\sigma_g}
  \right)
  \label{eq:mle_genuine}
\end{equation}
where $\phi(\cdot)$ is the standard normal density and $\hat{\alpha}_i$, $\hat{\lambda}_t$ are absorbed item and year fixed effects. In practice, this reduces to OLS of $\log b_{jt}$ on $\mathbf{x}_j$ with fixed effects, yielding $\hat{\sigma}_g^2$ as the residual variance.

\paragraph{Stage 2: Cover bid distribution.}
Conditional on the winning bid $b^*_t$, we estimate cover-bid parameters from FL bids under each regime specification (Equations~\ref{eq:cover_dist_r1}--\ref{eq:cover_dist_r2}).

\paragraph{Model selection.}
We select between regimes using BIC; the implied markup is:
\begin{equation}
  \hat{\mu}_{\text{structural}} =
  \frac{E[b^*_t \mid \text{FL present}] -
        E[b^*_t \mid \text{FL absent}]}{E[b^*_t \mid \text{FL absent}]}
  \label{eq:structural_markup}
\end{equation}
If the structural model is correctly specified, this markup should fall within the OLS--IV range, providing a structural consistency check. Standard errors are computed via parametric bootstrap (500 replications, tender-level resampling).

% ── 5.2 Reduced-Form Identification ──────────────────────────────

\subsection{Reduced-Form Identification}
\label{sec:reduced_form}

\paragraph{OLS baseline.}
Our baseline specification estimates the conditional association between
FL presence and tender outcomes:
\begin{equation}
  y_{igt} = \beta \cdot \text{losers}_{igt} + \mathbf{x}_{igt}'
  \boldsymbol{\delta} + \alpha_g + \lambda_t + \gamma_k + \varepsilon_{igt}
  \label{eq:ols_main}
\end{equation}
where $y_{igt}$ is the outcome for tender-item $i$ in item group $g$ at
time $t$ and purchasing unit $k$; $\text{losers}_{igt} \in \{0,1\}$
indicates FL presence; $\mathbf{x}_{igt}$ includes a convite indicator;
$\alpha_g$, $\lambda_t$, and $\gamma_k$ are item, year, and PBU fixed
effects; and $\varepsilon_{igt}$ is the error term clustered at the item
level.

The OLS coefficient $\hat{\beta}$ identifies a conditional
association---the average difference in outcomes between FL-present and
FL-absent tenders within the same item type, year, and purchasing unit.
Under Prediction~\ref{pred:price_v6}, $\hat{\beta} > 0$ for prices. The
OLS estimate of 0.064 (6.4\%) represents a lower bound on the causal
effect if the FL indicator contains measurement error (binary treatment
misclassification). We assess the sensitivity of $\hat{\beta}$ to
unobserved confounders using the \citet{cinelli2020making} framework,
which yields a robustness value of $RV_{q=1} = 17.5\%$: an unobserved
confounder would need to explain at least 17.5\% of the residual
variation in both FL presence and log prices to reduce $\hat{\beta}$ to
zero.

\paragraph{Matching estimators.}
To address potential selection on observables, we estimate the
FL--price effect using two matching approaches: coarsened exact
matching (CEM) on year, procedure type, item group, and PBU size
quartile ($\hat{\beta}_{\text{CEM}} = 0.077$, $N = 969{,}751$), and
inverse probability weighting (IPW, $\hat{\beta}_{\text{IPW}} = 0.055$,
$N = 830{,}194$). These estimates bracket the OLS baseline (0.064)
and the cross-fit average (0.036), providing a non-parametric check
that the FL--price association survives reweighting on observables.

\paragraph{Instrumental variable (supplementary).}
As a supplementary bracketing exercise, we instrument FL presence with
a leave-one-out measure of FL supply:
\begin{equation}
  Z_{kgt} = \sum_{j \neq k} \mathbf{1}[\text{FL firm active at PBU } j
  \text{ in group } g, \text{ year } t]
  \label{eq:iv_v6}
\end{equation}
The identifying assumption is that FL activity at other PBUs in the same
product market and year affects focal-PBU prices only through the
probability of FL participation---a supply-side exclusion restriction in
the spirit of \citet{goldsmithpinkham2020bartik}. We treat the IV
estimate as a Local Average Treatment Effect for complier tenders: an
upper bound on the average price effect that corrects for
measurement-error attenuation in the FL indicator.

\paragraph{IV validity concerns.}
Balance tests (Table~\ref{tab:iv_balance} in Appendix~\ref{app:robustness}) reveal systematic but economically negligible imbalance across all four observables (all standardized differences $< 0.03\sigma$). The \citet{goldsmithpinkham2020bartik} Rotemberg decomposition is infeasible in our setting ($> 200{,}000$ item-year shares). We therefore rely on matching (CEM, IPW) and sensitivity analysis (\citealt{cinelli2020making}, $RV = 17.5\%$) to bound IV validity. The IV is supplementary; our primary estimates are OLS (0.064), cross-fit (0.036), and matching (0.055--0.077).

% ── 5.3 Detection Validation ─────────────────────────────────────

\subsection{Detection Performance Validation}
\label{sec:detection_validation}

We evaluate the FL screen as a detection tool using CADE cartel
convictions as ground truth. The validation uses the temporal split
described in Section~\ref{sec:cade_sample}: FL status is defined using
2009--2014 data and evaluated against 2015--2019 CADE outcomes.

\paragraph{ROC analysis.}
We vary the IQR multiplier from 0.5 to 5.0 in steps of 0.1, producing
46 different FL thresholds. For each threshold, we compute:
\begin{align}
  \text{TPR}(m) &= \frac{|\{j : \text{FL}_m(j)=1\} \cap
    \{j : \text{CADE}(j)=1\}|}{|\{j : \text{CADE}(j)=1\}|}
  \label{eq:tpr} \\
  \text{FPR}(m) &= \frac{|\{j : \text{FL}_m(j)=1\} \cap
    \{j : \text{CADE}(j)=0\}|}{|\{j : \text{CADE}(j)=0\}|}
  \label{eq:fpr}
\end{align}
where $\text{FL}_m(j)$ indicates FL classification under multiplier $m$
and $\text{CADE}(j)$ indicates co-participation with CADE-convicted
cartelists. AUC~$= 0.5$ corresponds to random classification and
AUC~$= 1.0$ to perfect discrimination.

\paragraph{Optimal threshold.}
We determine the welfare-maximizing IQR multiplier using Youden's $J$
statistic:
\begin{equation}
  m^*_{\text{optimal}} = \operatorname*{arg\,max}_m \;
  \left[\text{TPR}(m) - \text{FPR}(m)\right]
  \label{eq:youden}
\end{equation}
which balances true-positive and false-positive rates, connecting the
IQR rule to a formal detection-theoretic criterion.

\paragraph{Comparison with Imhof screens.}
We compare the FL screen with \citet{imhof2019detecting}-style
bid-level features (within-tender CV, kurtosis, skewness, normalized
spread) on the same CADE ground truth (Section~\ref{sec:results_detection}).

% ── 5.4 Bajari-Ye Tests ──────────────────────────────────────────

\subsection{Bajari--Ye Bid Coordination Tests}
\label{sec:bajari_ye}

We implement the \citet{bajari2003deciding} tests using corrected
first-stage residuals from the bid equation:
\begin{equation}
  \log b_{jt} = \mathbf{x}_j'\boldsymbol{\phi} + \alpha_i + \lambda_t
  + u_{jt}
  \label{eq:by_first_stage}
\end{equation}
where $\mathbf{x}_j$ includes firm size, firm age, and CNAE sector
dummies. We exclude $n_{\text{bids}}$ from the first stage: it is
jointly determined with bid levels and mechanically increased by FL
presence \citep[p.~978]{bajari2003deciding}.

Using the residuals $\hat{u}_{jt}$, we test (i)~\emph{exchangeability}
(KS test comparing FL and non-FL residual distributions;
Prediction~\ref{pred:exchange_v6}) and (ii)~\emph{conditional
independence} (mean pairwise product of FL residuals within tenders;
bootstrap with 1,000 iterations). A fake-groups placebo splits non-FL
firms randomly within each tender; the economically relevant statistic
is the FL--non-FL pairwise \emph{difference}, which isolates excess
FL correlation beyond common tender-level shocks. We also report a
tender-FE variant as a conservative specification absorbing all
tender-level variation.

% ── 5.5 Network-Split Heterogeneity ──────────────────────────────

\subsection{Network-Split Heterogeneity Test}
\label{sec:empirical_network}

To test Proposition~\ref{prop:market_selection}---that cover bidding
concentrates in competitive markets---we split FL firms into two
subgroups based on co-bidding network characteristics.

\paragraph{Winner HHI.}
For each FL firm $j$, we compute the Herfindahl--Hirschman Index of
winning firms across all tenders in which $j$ participates:
\begin{equation}
  \text{HHI}_j = \sum_{w=1}^{W_j} s_{wj}^2
  \label{eq:winner_hhi}
\end{equation}
where $s_{wj}$ is the share of firm $j$'s tenders won by winner $w$,
and $W_j$ is the number of distinct winners. A high $\text{HHI}_j$
indicates that $j$ operates in markets dominated by a single or few
winners (concentrated markets); a low $\text{HHI}_j$ indicates
diversified winner patterns (competitive markets).

\paragraph{Classification.}
We classify FL firms as \emph{concentrated-market} if
$\text{HHI}_j$ exceeds the sample median, and
\emph{competitive-market} otherwise. We then estimate the baseline
price regression (Equation~\ref{eq:ols_main}) separately for each
subgroup, interacting $\text{losers}_{igt}$ with the
market-structure indicator.
